\centerline{\Huge \bf  Тренировочная Олимпиада 2}
\centerline{\Huge \bf 8 класс}
\centerline{\Huge \bf Уровень: Регион ВсОШ. День 1}

\begin{flushright}
    \scriptsize{Глава 2. Зверь.}
\end{flushright}


\problem{1} 
 Докажите, что при $x, y \geq \dfrac{1}{2}$ выполняется неравенство:
 
\[ x^2y^2 + 2(x+y) \geq 4xy + 1 \]
\begin{flushright}
    \textit{\small{(Гасанов Р.В.)}}
\end{flushright}


\problem{2}
Остроугольный треугольник $ABC$ с центром описанной окружности в точке $O$ и центром вписанной окружности в точке $I$. Через точку $O$ проведена прямая, перпендикулярная $BI$, а через точку $I$ $-$ прямая, перпендикулярная $BO$. Оказалось, что они пересеклись в точке $A$. Найдите углы треугольника.
\begin{flushright}
    \textit{\small{(Гасанов Р.В.)}}
\end{flushright}


\problem{3} 
\textbf{Взлома кода Галактики.} Алина, подружившись с Гроксами попала в их огромную галактическую империю. Все планеты гроксов пронумерованы от 1 до 1000000000. Чтобы добраться от планеты под номером $a$ до планеты под номером $b$ нужно потратить НОД$(a, b)$ спорлингов. Может ли Алина, зная только цены перелётов между планетами, восстановить нумерацию?
\begin{flushright}
    \textit{\small{(Гасанов Р.В.)}}
\end{flushright}

\problem{4} 
Три квадратных трехчлена $P(x), Q(x), R(x)$ со старшими коэффициентами, равными 1, имеют по два действительных корня и удовлетворяют условиям 
\[P(Q(x)) = (x-1)(x-3)(x-5)(x-7), \ Q(R(x)) = (x-2)(x-4)(x-6)(x-8) \]

при всех действительных $x$. Чему может быть равно число $P(0)+Q(0)+R(0) $
\begin{flushright}
    \textit{\small{(Гасанов Р.В.)}}
\end{flushright}

\problem{5} 
Расул Владимирович составил каждому из 14 учеников индивидуальный вариант тренировочной Олимпиады, но при раздаче каким-то ученикам выдал не свои условия. Назовём обмен \textit{шкибидишным}, если некоторые (не обязательно все) ученики разибиваются на пары и внутри этой пары меняются между собой условиями. Докажите, что за два шкибидишных обмена можно сделать так, что у каждого будет его условие.
\begin{flushright}
    \textit{\small{(Гасанов Р.В.)}}
\end{flushright}